\section{Mathematical \& Technical Foundation}
\cvitem{}{
From the MSc in Computer Science and Engineering (Linköping University, 180 credits, Software Engineering
track), the following coursework directly underpins mathematics and physics teaching:
}
\cvitem{Mathematics}{
Calculus (one variable and several variables), Linear Algebra, Discrete Mathematics and Logic,
Probability, Statistics, Numerical Algorithms, Combinatorial Optimization, Transform Theory, Signal Theory
}
\cvitem{Electronics \& Physics}{
Basic Electronics and Measurement Technology, Circuit Theory, Switching Theory and Logical Design,
Computer Hardware and Architecture, Digital Project Laboratory, Microcomputer Project Laboratory,
Electromagnetism, Modern Physics, Automatic Control
}
\cvitem{Programming \& Applied Logic}{
Data Structures and Algorithms, Object-Oriented Development, Concurrent Programming and Operating
Systems, Computer Graphics, Artificial Intelligence, Database Technology, Computer Security
}
\cvitem{Master's Thesis}{
\emph{Evolving an OS Kernel Using Temporal Logic and Aspect-Oriented Programming} -- applying formal
(temporal) logic to software verification, at École des Mines de Nantes, France.
}
\cvitem{}{
Earlier, at upper-secondary level (Naturvetenskapligt program, Rönneskolan), coursework included
mathematics, physics and chemistry, alongside a technical elective in programming and electronics. A
special project built 3D/vector graphics in assembler on the Amiga, applying rotation matrices, line-drawing
algorithms and coordinate transformations -- linear algebra put directly into visual, hands-on practice.
}
